\begin{table}[H]
\centering
\caption[Multi-Seed ATE Stability (N=20)]{Height ATE distribution across 20 independent random seeds. Wide standard deviation reflects genuine effect heterogeneity, not estimator instability. Only the random seed varies; all other pipeline parameters are fixed.}
\label{tab:seed_stability}
\begin{tabular}{lcc}
\toprule
\textbf{Statistic} & \textbf{Height ATE (feet)} & \textbf{Sign Consistency} \\
\midrule
Mean (N=20) & $-$2.86  & 45\% positive \\
Std Dev     & \phantom{$-$}45.14 & --- \\
\bottomrule
\end{tabular}
\par\smallskip
\footnotesize Fixed pipeline specification. The wide standard deviation and near-symmetric sign consistency are expected when true causal effects are genuinely heterogeneous. As documented in Section~\ref{sec:hte}, $\hat{\tau}(x_i)$ ranges from $-$23.6 to 178.5 feet across the parcel distribution; a single aggregate ATE is poorly identified because marginal sampling variation shifts the aggregate substantially. The Causal Forest is designed to estimate the function $\hat{\tau}(x_i)$, not a scalar point treatment effect. The correct summary is the CATE quartile distribution in Table~\ref{tab:cate_distribution}, where the median is stably 16.76 feet and 91.8\% of parcels show a consistent direction.
\end{table}
